% Generated by `taskgen course examples` from dlp.shanks_subgroup (seed 3).
% Do not edit: change examples.yaml and regenerate.
\begin{example}[алгоритм Шенкса в подгруппе]\label{ex:bsgs-subgroup}
Известно, что $\ord(38) = 83$ в $(\mathbb{Z}/167\mathbb{Z})^*$. Решить методом Шенкса \[38^x \equiv 120\ (\textrm{mod}\ 167)\] (или показать, что решений нет), выбрав длину списков по порядку $g$, а не по $p - 1$.

\par\noindent\textit{Решение.} Решение ищем среди $0 \le x < n = 83$; $m = \lceil\sqrt{83}\rceil = 10$ (вместо $\lceil\sqrt{166}\rceil = 13$), так как каждое $x < n \le m^2$ записывается как $x = im + j$, $0 \le i, j < m$. Малые шаги $g^j$:\[\begin{array}{c|rrrrrrrrrr}j & 0 & 1 & 2 & 3 & 4 & 5 & 6 & 7 & 8 & 9 \\ \hline g^j & 1 & 38 & 108 & 96 & 141 & 14 & 31 & 9 & 8 & 137\end{array}\]$g^{-1} \equiv 22$, $g^{-m} \equiv 22^{10} \equiv 144$. Большие шаги $y_i = h \cdot g^{-im}$ ($y_{i+1} = y_i \cdot 144$):\[\begin{array}{c|rrrrrrrrrr}i & 0 & 1 & 2 & 3 & 4 & 5 & 6 & 7 & 8 & 9 \\ \hline y_i & 120 & 79 & 20 & 41 & 59 & 146 & 149 & 80 & 164 & 69\end{array}\]Все $m$ больших шагов сделаны, совпадений нет. Если бы $h = g^x$, $0 \le x < n$, то при $x = im + j$ было бы $y_i = g^j$; значит, $h \notin \langle g \rangle$. Независимая проверка: $\langle g \rangle$ --- единственная подгруппа порядка $n$ циклической группы $(\mathbb{Z}/167\mathbb{Z})^*$, она состоит из решений $y^{n} = 1$, а $120^{83} \equiv 166 \ne 1$.\[\text{Решений нет.}\]
\end{example}
